\chapter{Introduction}

Quantum computing is a promising paradigm that offers potentially exponential advantages over classical computing for problems with a specific structure, such as factorization and quantum simulation. A key property enabling these benefits is the linearity of quantum mechanics, which, in principle, prevents errors from accumulating between gates \citep{Kitaev2002}. However, the inherent quantum nature and extreme sensitivity of these systems to errors—introduced by gates or external conditions such as environmental fluctuations and electromagnetic interference—mean that linearity alone is insufficient to ensure reliable results, particularly for large-scale operations like factorization \citep{Shor1996}. Furthermore, the no-cloning theorem prohibits the copying of qubits, rendering traditional error-correction methods, such as redundancy, inapplicable in a straight forward way.

Quantum error-correction is a field developed to address this challenge. It has been demonstrated that quantum information can be encoded non-locally by distributing it across multiple qubits \citep{Shor1996}. Codes that employ this strategy are known as quantum error-correcting codes. The efficacy of this method is supported by recent experimental improvements in physical quantum computers \citep{IBM2025}. Nevertheless, it remains an open area of research, and advancements in this field are critical for the continued development of quantum computing \citep{IBM2025}.

This thesis presents the theoretical foundations of quantum error-correcting codes as a key component for fault-tolerant quantum computing. Rather than focusing on hardware implementations, this work explores the algebraic structure and mathematical basis of the codes. It starts from the error models and fundamental limitations imposed by the underlying quantum system, culminating in the efficient formulation of the stabilizer formalism and its corresponding codes.

To this end, this thesis is structured as follows. Chapter 2 defines the error model that will permeate the subsequent chapters. Chapter 3 establishes the mathematical basis of the field, defining what a code is and the conditions it must fulfill to correct errors. Chapter 4 introduces linear codes, a classical concept that serves as inspiration for some of the most important quantum error-correcting codes. Chapter 5 defines the stabilizer formalism and stabilizer codes, a generalization of linear codes into quantum computing. Chapter 6 dives into the first code example, the repetition code, from the classical system to the quantum circuit that implements it. Chapter 7 develops the Shor code as a prime example of both a general and a stabilizer code. Lastly, Chapter 8 presents the conclusions.
